\section{A ternary line-code obstruction}\label{sec:line-code-obstruction}

We now sharpen Proposition~\ref{prop:characteristic-three-modular-bound}.
The preceding repair argument approaches the center
\[
 k_0(q)=\frac{q^2}{3}+\frac{4q}{3}.
\]
At this center, the first two secant moments produce a small-weight word in
the ternary code generated by the lines of \(\PG(2,q)\).  Two pointwise
integer-shell inequalities force that word to need at least three generator
lines and show that three lines already cost an additional \(q/3-o(q)\)
points.

Fix \(q=3r\), put \(n=2r+1\), and let \(\mathcal A\) be a complete
\((K,n)\)-arc.  Let \(\mathcal D\) be the point set in the dual plane
corresponding to the full family of \(n\)-secants.  After fixing a polarity,
write \(M\) for the symmetric point--line incidence matrix and
\[
 d=M\mathbf 1_{\mathcal D},\qquad T=|\mathcal D|.
\]
Thus \(d_P\) is the number of \(n\)-secants through \(P\), and completeness
gives \(d_P\ge1\) for \(P\notin\mathcal A\).  The incidence and pair
identities are
\begin{equation}\label{eq:line-code-degree-moments}
 \sum_{P\in\mathcal A}d_P=nT,\qquad
 \sum_{P\notin\mathcal A}d_P=rT,\qquad
 \sum_P\binom{d_P}{2}=\binom T2.
\end{equation}

Write
\begin{equation}\label{eq:line-code-offsets}
 K=3r^2+4r+\delta,\qquad T=6r+j,\qquad
 u=\mathbf1+3\mathbf1_{\mathcal A}-d.
\end{equation}

\begin{lemma}[Centered moment identities]\label{lem:centered-moment-identities}
\coverage{fragment}
With the notation above,
\begin{align}
 \sum_Pu_P&=(9-3j)r+3\delta-j+1,
 \label{eq:centered-sum}\\
 \sum_Pu_P^2&=(15-3j)r+15\delta+j^2-8j+1,
 \label{eq:centered-norm}\\
 \sum_{P\in\mathcal A}u_P&=(10-2j)r+4\delta-j.
 \label{eq:centered-internal-sum}
\end{align}
Moreover, the shell defect
\begin{equation}\label{eq:shell-defect-definition}
 E=\sum_{P\in\mathcal A}\frac{u_P(u_P-1)}2
   +\sum_{P\notin\mathcal A}\frac{u_P(u_P+1)}2
\end{equation}
is nonnegative and equals
\begin{equation}\label{eq:shell-defect}
 E=(2-j)r+1+5\delta+\frac{j(j-7)}2.
\end{equation}
\lean{centeredMoment_sum,centeredMoment_norm,centeredMoment_internalSum,
shellDefect_identity}
\end{lemma}

\begin{proof}
The first and third identities follow from
\(\sum_Pd_P=(q+1)T\) and the first identity in
\eqref{eq:line-code-degree-moments}.  The pair identity gives
\(\sum_Pd_P^2=T(T+q)\); expanding \(u^2\) gives
\eqref{eq:centered-norm}.  Expanding
\eqref{eq:shell-defect-definition} and applying the same moments gives
\eqref{eq:shell-defect}.  Every shell summand is nonnegative: on the
complement, completeness gives \(u_P=1-d_P\le0\).
\end{proof}

Let \(C_1(2,q)\) denote the ternary code generated by the line incidence
vectors.  Reducing \(u\) modulo three gives
\begin{equation}\label{eq:residue-word}
 \bar u=\mathbf1-M\mathbf1_{\mathcal D}\in C_1(2,q),
\end{equation}
because the sum of the lines in a pencil is \(\mathbf1\) modulo three.
The identity \(M^2=qI+J\) also gives
\begin{equation}\label{eq:incidence-of-centered-word}
 Mu=(q+1-T)\mathbf1+3M\mathbf1_{\mathcal A}-q\mathbf1_{\mathcal D}.
\end{equation}

We use the following exact form of the small-codeword theorem.  If a nonzero
word \(c\in C_1(2,q)\), where \(q=p^e>27\) and \(e>2\), has weight
\[
 w(c)<(\lfloor\sqrt q\rfloor+1)
       (q+1-\lfloor\sqrt q\rfloor),
\]
then Sz\H{o}nyi--Weiner~\cite[Theorem~4.3]{SWmodp} express \(c\) as a linear
combination of exactly \(\lceil w(c)/(q+1)\rceil\) distinct lines.  In
characteristic three its nonzero coefficients have balanced representatives
in \(\{-1,1\}\).

\begin{theorem}[Three-line gap]\label{thm:line-code-gap}
\coverage{fragment}
Fix \(C>0\) and \(\varepsilon>0\).  For all sufficiently large
\(q=3^e\), there is no complete \((K,2q/3+1)\)-arc in \(\PG(2,q)\) with
\begin{equation}\label{eq:line-code-gap}
 K=\frac{q^2}{3}+\frac{4q}{3}+\delta,\qquad
 -Cq\le\delta\le\left(\frac13-\varepsilon\right)q.
\end{equation}
\imports{SWmodp:small-codeword}
\lean{threeLine_shellCollapse,threeLine_phaseBoundary}
\end{theorem}

\begin{proof}
Suppose otherwise, and write \(\delta=\alpha q\).  Applying
Cauchy--Schwarz separately to the two classes in
\eqref{eq:line-code-degree-moments} gives
\begin{equation}\label{eq:line-code-cauchy}
 F(T):=\frac{n^2T^2}{K}-nT
 +\frac{r^2T^2}{q^2+q+1-K}-rT-T(T-1)\le0.
\end{equation}
Coverage gives \(T\ge(q^2+q+1-K)/r\), so it bounds
\(j=T-6r\) below in terms of \(C\).  Write \(F(T)=AT^2-3rT\), with
\[
 A=\frac{n^2}{K}+\frac{r^2}{q^2+q+1-K}-1,
\]
which is positive by the same weighted Cauchy argument as in
Section~\ref{sec:n-secants}.  Since \(F(T)\le0\) and \(T>0\),
\(T\le3r/A\).  Uniformly for \(\alpha\) in the compact interval prescribed
by \eqref{eq:line-code-gap}, direct expansion gives
\[
 \frac{3r}{A}=6r+5+15\alpha+O_{C,\varepsilon}(r^{-1}).
\]
Thus \(j\) is bounded above as well.  Pass to a subsequence on which \(j\)
is fixed.  For later reference, direct expansion now gives
\[
 F(6r+j)=3(j-5-15\alpha)r+O_{C,\varepsilon,j}(1).
\]
The centered identities then give
\(E=O(q)\) and \(\sum_Pu_P^2=O(q)\).

Let \(z_P\in\{-1,0,1\}\) be the balanced representative of \(u_P\pmod3\),
put \(R=\{P:z_P\ne0\}\), and write \(w=|R|\).  For the summand \(e_P\) of
\eqref{eq:shell-defect-definition}, the pointwise inequality
\begin{equation}\label{eq:shell-support-lower}
 2\mathbf1_{\mathcal A}(P)u_P-e_P-u_P
 \le\mathbf1_R(P)
\end{equation}
holds.  On the arc its left side is
\(u_P-u_P(u_P-1)/2\), while off the arc it is
\(-u_P(u_P+3)/2\); the three residue classes prove the claim.  Summing and
using Lemma~\ref{lem:centered-moment-identities} gives the exact cancellation
\begin{equation}\label{eq:three-line-support}
 3q-\frac{(j-1)(j-4)}2
 =2\sum_{P\in\mathcal A}u_P-E-\sum_Pu_P\le w.
\end{equation}

By \eqref{eq:residue-word}, \(\bar u\) belongs to the ternary line code, and
\(w\le\sum_Pu_P^2=O(q)\).  Sz\H{o}nyi--Weiner's theorem therefore gives
\[
 \bar u=\sum_{i=1}^t\varepsilon_i\mathbf1_{L_i},\qquad
 t=\left\lceil\frac{w}{q+1}\right\rceil,\qquad
 \varepsilon_i\in\{-1,1\},
\]
on distinct lines \(L_i\).  The norm bound makes \(t\) bounded.  Away from
their \(O(t^2)\) intersection points, every generator line contributes a
nonzero coordinate, so
\begin{equation}\label{eq:line-code-weight}
 w=tq+O(1).
\end{equation}
Equation~\eqref{eq:three-line-support} forces
\begin{equation}\label{eq:three-lines-needed}
 t\ge3.
\end{equation}

Put \(\sigma=\sum_i\varepsilon_i\), and let
\(a=(t+\sigma)/2\) be the number of positive generator lines.  Pointwise,
\begin{equation}\label{eq:signed-norm-correction}
 (u_P^2-z_P^2)-(u_P-z_P)\ge0.
\end{equation}
Indeed \(u_P\equiv z_P\pmod3\); if \(z_P=1\) and
\(P\notin\mathcal A\), completeness gives \(u_P\le-2\), and the left side
is at least six.  The positive generator lines contain \(aq+O(1)\)
nonintersection points with \(z=1\), but at most \(an\) of them lie on the
arc.  Summing \eqref{eq:signed-norm-correction} consequently gives
\begin{equation}\label{eq:positive-line-cost}
 \left(\sum_Pu_P^2-w\right)-\left(\sum_Pu_P-\sum_Pz_P\right)
 \ge2aq-O(1).
\end{equation}
Since \(\sum_Pz_P=\sigma q+O(1)\), substitution of
\eqref{eq:centered-sum}, \eqref{eq:centered-norm}, and
\eqref{eq:line-code-weight} into \eqref{eq:positive-line-cost} cancels
\(j\) and \(\sigma\), leaving
\begin{equation}\label{eq:three-line-upper}
 t\le1+6\alpha+o(1).
\end{equation}
Together with \eqref{eq:three-lines-needed}, this gives
\(\alpha\ge1/3-o(1)\), contrary to
\(\alpha\le1/3-\varepsilon\).
\imports{SWmodp:small-codeword}
\uses{lem:centered-moment-identities,lem:pencil-bounds}
\end{proof}

\begin{proposition}[Sublinear-repair rigidity]
\label{prop:three-line-rigidity}
\coverage{fragment}
Suppose that a sequence of complete
\((K,2q/3+1)\)-arcs over \(q=3^e\) satisfies
\[
 K=\frac{q^2}{3}+\frac{5q}{3}+o(q).
\]
Then, for every sufficiently large member of the sequence, the centered
residue word \(\bar u\) is an exact signed combination of three distinct
lines.
\imports{SWmodp:small-codeword}
\lean{threeLine_rigidityBoundary}
\end{proposition}

\begin{proof}
Coverage and \eqref{eq:line-code-cauchy} again keep \(j=T-2q\) bounded.
Equations~\eqref{eq:three-lines-needed} and
\eqref{eq:three-line-upper}, now with \(\alpha=1/3+o(1)\), give
\(3\le t\le3+o(1)\).  Since \(t\) is an integer, it equals three for all
sufficiently large fields.  The exact small-codeword theorem supplies the
claimed signed representation.
\end{proof}

\begin{proposition}[Exact endpoint obstruction]
\label{prop:exact-five-thirds-endpoint}
\coverage{absent}
For all sufficiently large \(q=3^e\), there is no complete
\[
 \left(\frac{q^2}{3}+\frac{5q}{3},\frac{2q}{3}+1\right)\text{-arc}
 \quad\text{in }\PG(2,q).
\]
\end{proposition}

\begin{proof}
Suppose that such arcs exist for arbitrarily large \(q\), and pass to a
sequence of them.  Put \(q=3r\) and \(\delta=r\) in
\eqref{eq:line-code-offsets}.  Proposition~\ref{prop:three-line-rigidity}
shows that, for every sufficiently large member of the sequence,
\[
 z=\varepsilon_1\mathbf1_{L_1}
   +\varepsilon_2\mathbf1_{L_2}
   +\varepsilon_3\mathbf1_{L_3}\pmod3,
 \qquad \varepsilon_i\in\{-1,1\},
\]
where \(z_P\in\{-1,0,1\}\) is the balanced representative of \(u_P\).
Write \(\sigma=\varepsilon_1+\varepsilon_2+\varepsilon_3\).

We first record the exact scalar reduction.  Put
\(R_+=\{P:z_P=1\}\), \(R_-=\{P:z_P=-1\}\), and let
\(a=(3+\sigma)/2\) be the number of positive generator lines.  Pointwise
on \(\mathcal A\), one has
\(u_P\le e_P+\mathbf1_{R_+\cup R_-}(P)\); off \(\mathcal A\), one has
\(e_P\ge1\) on \(R_+\).  Therefore
\begin{equation}\label{eq:endpoint-shell-cap}
 \sum_{P\in\mathcal A}u_P
 \le E-|R_+|+2|\mathcal A\cap R_+|
       +|\mathcal A\cap R_-|.
\end{equation}
The signed three-line representation gives
\(|R_+|=aq+O(1)\).  The arc cap on its three generator lines gives
\[
 |\mathcal A\cap R_+|\le a(2q/3+1)+O(1),\qquad
 |\mathcal A\cap R_-|\le(3-a)(2q/3+1)+O(1).
\]
Substitution of the exact centered identities in
\eqref{eq:endpoint-shell-cap} yields \(j\ge1+a\) for all sufficiently
large \(q\).

Also \(E\ge0\) gives \(j\le7\), while summing the residue word and comparing
\(u\) with its balanced representative give
\begin{equation}\label{eq:endpoint-scalar-conditions}
 \sigma\equiv1-j\pmod3,
 \qquad
 |4-j-\sigma|\le7-j.
\end{equation}
Indeed the second inequality is the coefficient of \(q\) in
\(\left|\sum_P(u_P-z_P)\right|\le\sum_P(u_P^2-z_P^2)\).
Checking the four values \(\sigma\in\{-3,-1,1,3\}\) under
\(1+a\le j\le7\) and \eqref{eq:endpoint-scalar-conditions} leaves
\[
 (j,\sigma)=(1,-3),(2,-1),(3,1),(4,-3),(4,3),(5,-1),(7,-3).
\]
For the last row, \(\sum_Pu_P^2=3q-6<w\), so the six possibilities are
\begin{equation}\label{eq:endpoint-six-rows}
 (j,\sigma)=(1,-3),(2,-1),(3,1),(4,-3),(4,3),(5,-1).
\end{equation}

The three lines cannot be concurrent.  When their signs agree, the exact
inequality above gives the following contradictions; the two numerical
columns are \(\left|\sum u-\sum z\right|\) and
\(\sum u^2-w\), respectively:
\[
\begin{array}{c|c|c}
 (j,\sigma)&\left|\sum u-\sum z\right|&\sum u^2-w\\ \hline
 (1,-3)&6q&6q-6\\
 (4,-3)&3q-3&3q-15\\
 (4,3)&3q+3&3q-15.
\end{array}
\]
For the mixed-sign rows, the arc cap leaves at least \(q/3-1\) external
positive nonvertices when \(j=2,5\), and at least \(2q/3-1\) external
positive residues when \(j=3\).  Substitution in the exact pointwise
correction \eqref{eq:signed-norm-correction} makes
\eqref{eq:positive-line-cost} fail by six in each case.

The generators therefore form a triangle.  Its residue support has size
\(3q\) when the signs agree and \(3q-2\) when they are mixed.  Comparing
with \eqref{eq:three-line-support} immediately excludes the rows
\((2,-1)\) and \((3,1)\).  In each remaining row equality holds in the
pointwise inequality \eqref{eq:shell-support-lower}.  Checking its three
residue classes gives the exact degree rule
\begin{equation}\label{eq:endpoint-equality-degrees}
 z_P=1\Longrightarrow d_P=3,\qquad
 z_P=-1\Longrightarrow d_P=2,\qquad
 z_P=0\Longrightarrow d_P\in\{1,4\}.
\end{equation}

Dualize the generator lines to noncollinear points \(Q_1,Q_2,Q_3\), put
\(p_i=(1+\varepsilon_i)/2\), and put
\(y_i=\mathbf1_{\mathcal D}(Q_i)\).  If \(x_{ij}\) is the
\(\mathcal D\)-degree of the connector \(Q_iQ_j\), then the \(q-1\)
ordinary lines through \(Q_i\) have degree \(2+p_i\) by
\eqref{eq:endpoint-equality-degrees}.  Counting incidences with
\(\mathcal D\) in that pencil gives the general identity
\begin{equation}\label{eq:endpoint-pencil-equation}
 x_{ij}+x_{ik}=j+2+p_i+q(y_i-p_i),
 \qquad \{i,j,k\}=\{1,2,3\}.
\end{equation}

For \((j,\sigma)=(1,-3)\), all three connectors have degree three, so
\eqref{eq:endpoint-pencil-equation} would give \(6=3+qy_i\), impossible
for \(q>3\).  For \((4,3)\), all connectors have degree two, and the same
equation gives \(4=7+q(y_i-1)\), again impossible for \(q>3\).
For \((4,-3)\), it gives \(y_i=0\) and connector degrees
\((3,3,3)\).  Finally take \((5,-1)\), with one positive generator.  At
its point \(Q_i\), the two connectors have degrees in \(\{1,4\}\), and
\eqref{eq:endpoint-pencil-equation} forces \(y_i=1\) and their sum to be
eight; both therefore have degree four.  At either negative generator the
remaining connector has degree three, so the same equation forces \(y_i=0\).
Thus the only surviving rows and connector profiles are
\[
 (j,\sigma)=(4,-3),(5,-1),\qquad
 (x_{12},x_{13},x_{23})=(3,3,3),(4,4,3),
\]
up to relabeling.

For \((j,\sigma)=(4,-3)\), the positive residue consists only of the three
triangle vertices.  On \(\mathcal A\), the pointwise inequality
\(u_P\le2e_P+\mathbf1_{\{z_P=1\}}\) therefore gives
\[
 2q-4=\sum_{P\in\mathcal A}u_P\le2E+3=2q-7,
\]
a contradiction.

It remains to exclude \((j,\sigma)=(5,-1)\).  Let \(N_i\) be the number of
lines of \(\mathcal D\)-degree \(i\).  The equality rule
\eqref{eq:endpoint-equality-degrees} gives
\(N_2=2q-2\) and \(N_3=q\), while the total-line and first-moment identities
give
\[
 N_4=\frac{q(q+2)}3+2.
\]
Let \(A_i\) count points of \(\mathcal A\) whose dual lines have
\(\mathcal D\)-degree \(i\), put \(b=A_1\), and put \(h=N_4-A_4\).
The selected-point and selected-incidence equations reduce to
\[
 A_3=\frac{2q}{3}+1+b+2h.
\]
Only one degree-three point lies off the positive generator line, so its arc
cap gives \(b+2h\le1\).  Hence \(h=0\), and both degree-four connector
vertices are selected.  At least \(2q/3\) selected degree-three points then
lie on the positive generator; the two connector vertices raise its
intersection with \(\mathcal A\) to at least \(2q/3+2\), contradicting the
maximum \(2q/3+1\).  This excludes the last row.
\uses{lem:centered-moment-identities,prop:three-line-rigidity}
\end{proof}

\begin{proof}[Proof of Theorem~\ref{thm:characteristic-three}]
Fix \(\varepsilon>0\).  If \(\varepsilon>1/3\),
Proposition~\ref{prop:characteristic-three-modular-bound} already gives the
claim, so assume \(0<\varepsilon\le1/3\).  If the theorem failed along an
unbounded sequence, that proposition would give
\[
 K\ge \frac{q^2}{3}+\frac{4q}{3}-o(q),
\]
whereas failure would give
\[
 K\le \frac{q^2}{3}+\frac{4q}{3}
       +\left(\frac13-\varepsilon\right)q.
\]
For large \(q\), these inequalities put \(\delta\) in the range of
Theorem~\ref{thm:line-code-gap}, with any fixed \(C>0\), a contradiction.
\imports{SWmodp:repair,SWmodp:small-codeword}
\uses{prop:characteristic-three-modular-bound,thm:line-code-gap}
\end{proof}
